% Generated draft: requires mathematical and source review.
\documentclass[11pt]{article}
\usepackage[margin=25mm]{geometry}
\usepackage{fontspec}
\setmainfont{Latin Modern Roman}
\usepackage{amsmath,amssymb,mathtools,unicode-math}
\setmathfont{Latin Modern Math}
\usepackage{enumitem,needspace,xurl,hyperref}
\setlength{\parindent}{0pt}
\setlength{\parskip}{5pt}
\setlength{\emergencystretch}{3em}
\newcommand{\N}{\mathbb{N}}
\newcommand{\Z}{\mathbb{Z}}
\newcommand{\Q}{\mathbb{Q}}
\newcommand{\R}{\mathbb{R}}
\newcommand{\C}{\mathbb{C}}
\newcommand{\F}{\mathbb{F}}
\newcommand{\A}{\mathbb{A}}
\newcommand{\PP}{\mathbb{P}}
\newcommand{\E}{\mathbb{E}}
\newcommand{\eps}{\varepsilon}
\newcommand{\dd}{\,\mathrm{d}}
\begin{document}
\section*{UnsolvedMath ID: 20003106}
\textbf{Source problem code:} catalogue entry
\subsection{Definitions and mathematical statement}
Fix a finite-type hyperbolic surface \(S\), with its Teichm\"uller space and the theorem range specified. Use the unscaled length-spectrum metric
\[
d_{\rm ls}(X,Y)=\max\{d_{\rm Th}(X,Y),d_{\rm Th}(Y,X)\}
=\sup_\gamma\left|\log\frac{\ell_\gamma(Y)}{\ell_\gamma(X)}\right|.
\]
A convention multiplying this by \(1/2\) gives the same isometry group. The question attributed to A. Papadopoulos is whether every bijective length-spectrum isometry comes from a mapping class. The intended candidate must be the effective extended group
\[
\operatorname{Mod}^{\pm}(S)/\ker\bigl(\operatorname{Mod}^{\pm}(S)\curvearrowright\mathcal T(S)\bigr),
\]
including orientation reversals and dividing out ineffective low-complexity kernels. Where orientation-reversing actions are effective, restricting to orientation-preserving mapping classes makes the proposed equality false. Degenerate and exceptional surface types require their own conventions; the source gives no surface range.

The supplied aid reports that every effective extended mapping class is a length-spectrum isometry, and that the equality with this corrected candidate remains open in the literature it checked. It records the following separation. The subgroup of maps preserving the directed distance \(d_{\rm Th}\) globally is the effective extended mapping class group, in the applicable Thurston-rigidity range. Maps reversing that distance globally satisfy
\[
d_{\rm Th}(FX,FY)=d_{\rm Th}(Y,X)
\]
and form at most one coset of the forward-preserving subgroup, if any exist. A remaining length-spectrum isometry would be genuinely mixed: neither globally forward-preserving nor globally reversing. The existence or exclusion of that reversing coset is a separate issue.

The aid gives a conditional infinitesimal reduction for \(C^1\) isometries. Let
\[
p_X(v)=\sup_\gamma\frac{\dd\ell_\gamma|_X(v)}{\ell_\gamma(X)},\qquad
N_X(v)=\max\{p_X(v),p_X(-v)\}.
\]
Suppose the unordered pair of opposite gauges \(\{p_X,p_X\circ(-I)\}\) can be recovered from the symmetric norm \(N_X\) so that every derivative isometry of \(N\) either preserves or exchanges the two gauges, with a locally unique sign choice. Then for a \(C^1\) length-spectrum isometry \(F\), at every \(X\) either
\[
p_{FX}(D_XF(v))=p_X(v)\quad\forall v,
\]
or the same equality holds with \(p_X(-v)\) on the right. Under the stated local uniqueness/continuity hypothesis, connectedness of Teichm\"uller space forces one global sign. Integrating the gauge along paths then makes \(F\) either a forward Thurston isometry, hence a mapping-class action in the reported rigidity range, or a global forward-to-reverse anti-isometry. This does not prove gauge recoverability, \(C^1\) regularity of every metric isometry, or nonexistence of reversing maps. None may be omitted when trying to deduce the original full isometry-group equality.
\subsection{Short English statement}
Determine whether length-spectrum isometries are exactly effective extended mapping classes. The supplied reduction separates directed-distance-preserving, reversing and mixed maps; eliminating mixed maps requires a conditional recovery of the two asymmetric tangent gauges, and does not itself eliminate global reversals.
\subsection{Sources}
\begin{itemize}
\item A. Papadopoulos, Problem 3.3, AIM workshop \emph{Lipschitz metric on Teichmueller space}, \url{https://aimath.org/pastworkshops/lipschitzteichproblems.pdf}. Decorative dots and the following page/author header are not part of the question.
\item Local UnsolvedMath registry, permanent ID 20003106, exact code \texttt{catalogue entry}; canonical location \texttt{aim-topology-notes.json}, notes[193], dated 14 August 2026; statement labeled exact, verification identifies the source index and OCR wording.
\item Research attempt 1: reduction, problem open at the run, rigor label proved, novelty confidence low; registry status partially solved. The directed-rigidity account and sign-splitting reduction are machine-generated research aids, not peer reviewed; independent verification is required before citation. No full external bibliography or complete exceptional-surface list is supplied.
\end{itemize}
\end{document}
